\documentclass{article}
\usepackage{amsmath,amssymb}
\usepackage{xurl}
\usepackage[hidelinks]{hyperref}
% Shared commands for notes in docs/paper-gaps/.

\DeclareUnicodeCharacter{2080}{\ifmmode{}_{0}\else\textsubscript{0}\fi}
\DeclareUnicodeCharacter{2081}{\ifmmode{}_{1}\else\textsubscript{1}\fi}
\DeclareUnicodeCharacter{2082}{\ifmmode{}_{2}\else\textsubscript{2}\fi}
\DeclareUnicodeCharacter{2099}{\ifmmode{}_{n}\else\textsubscript{n}\fi}

\newcommand{\Lean}{\textsc{Lean}}
\ExplSyntaxOn
\NewDocumentCommand{\leanid}{m}{
  \ifmmode
    \textsf{\detokenize{#1}}
  \else
    \group_begin:
    \urlstyle{sf}
    \tl_set:Nn \l_tmpa_tl {#1}
    \tl_replace_all:Nnn \l_tmpa_tl {\_} {_}
    \exp_args:NV \path \l_tmpa_tl
    \group_end:
  \fi
}
\ExplSyntaxOff
\providecommand{\tr}{\operatorname{tr}}
\newcommand{\tnleanIssueBase}{https://github.com/LionSR/TNLean/issues/}
\newcommand{\tnleanPrBase}{https://github.com/LionSR/TNLean/pull/}
\newcommand{\tnleanDocBase}{https://sirui-lu.com/TNLean/docs/}
\newcommand{\ghissue}[1]{\href{\tnleanIssueBase#1}{GitHub issue \##1}}
\newcommand{\ghpr}[1]{\href{\tnleanPrBase#1}{PR \##1}}
\newcommand{\doclink}[1]{\href{\tnleanDocBase#1.html}{\path{#1}}}

% Dirac notation and the identity operator, as in blueprint/src/macros/common.tex.
% \providecommand keeps note-local definitions authoritative.
\usepackage{dsfont}
\providecommand{\ket}[1]{|#1\rangle}
\providecommand{\bra}[1]{\langle #1|}
\providecommand{\braket}[2]{\langle #1 | #2 \rangle}
\providecommand{\ketbra}[2]{|#1\rangle\!\langle #2|}
\providecommand{\Id}{\mathds{1}}
\providecommand{\one}{\mathds{1}}
\providecommand{\C}{\mathbb{C}}
\providecommand{\MN}[1]{M_{#1}(\C)}
\providecommand{\MD}{M_D(\C)}

% Verdict marker: \gapnote{<kind>}{<status>}. Kind is the policy
% classification (clarification, local-correction, scope-restriction,
% unfaithful, false-source, open-gap); status is open, wip (elimination
% underway), resolved, or historical. Machine-read by the site index;
% prints a verdict line.
\newcommand{\gapnote}[2]{\par\noindent\textbf{Verdict:} #1 (#2).\par}

\title{Quantifiers in the Frustration-Free Ground-State Bond-Dimension Dichotomy}
\author{}
\date{2026-10-29}
\begin{document}
\maketitle
\gapnote{scope-restriction}{open}

\section{Source and quantified asymptotic result}

Theorem~5 of \cite{Sanz2010Quantum}, lines~850--880 of
\path{Papers/0909.5347/main.tex}, asserts a length-independent matrix size or
an order-$N^{1/5}$ obstruction at prime lengths for the unique ground states of
one fixed frustration-free local interaction.

The source-written quantified version takes a positive $L$-site interaction
$h$ and assumes that the zero-energy space of $H_M=\sum_i\tau^i(h)$ has
dimension one at every permitted volume $M\ge L$. It proves the alternative
\begin{align}
     & \exists B,M_0\quad
    \ker H_M=\C V^{(M)}(B),\quad V^{(M)}(B)\ne0
    \quad\text{for every }M\ge M_0; \label{eq:spwc10_eventual_uniform_branch} \\
     & \text{or}\quad
    N<L+\max\{q(D),3(D-1)(q(D)+1)\},\qquad
    q(D)=(D^2+1)^2, \label{eq:spwc10_prime_obstruction}
\end{align}
where (\ref{eq:spwc10_prime_obstruction}) quantifies over every prime $N$ and
every bond-$D$ tensor whose nonzero periodic vector is in $\ker H_N$.
The interaction and its range are fixed before those quantifiers. The
polynomial bound is an explicit fifth-degree obstruction; it is not a
statement only about the W state or a chosen sequence of representations.

\section{Derived data and the remaining boundary}

When the threshold in (\ref{eq:spwc10_prime_obstruction}) fails, primality and
$D<N$ remove every nonprimitive periodic block at length $N$. Grouping phase
classes and discarding zero combined coefficients produces an actual
expansion into separated blocks, with common injectivity length $q(D)$.
Complementary simultaneous spanning then gives $G_L(B_j)\subseteq\ker h$.
Consequently each block gives zero energy at every permitted volume. Its
injectivity ensures a nonzero periodic vector for $M\ge q(D)+1$, and the
one-dimensional ground-space hypothesis identifies the entire kernel with
that line. This derives (\ref{eq:spwc10_eventual_uniform_branch}) with
$M_0=\max\{L,q(D)+1\}$ and a block of bond dimension at most $D$.

The fixed-tensor conclusion is explicitly eventual. It does not identify the
finitely many shorter ground spaces with the same tensor, nor supply their
periodic trace-MPS representations at one common padded bond dimension.
A literal all-volume uniform-dimension formulation needs that finite-prefix
argument separately. The asymptotic alternative and its exact quantifiers
are therefore stated separately from an unrestricted source-labelled theorem.
No original canonical, phase-separation, transfer-gap, local-kernel or
future-volume equation is assumed.

The proof scripts and their Blueprint entries are pending hosted acceptance;
no local Lean build or elaboration was run. The full source theorem remains
open in issue \#10280 until the scope decision, validation and integration are
complete.

\section{Verdict}

\textbf{Verdict: scope restriction.} The source-written result gives the
quantified asymptotic dichotomy and an explicit fifth-degree obstruction.
The all-volume finite-prefix completion is not part of that result.

\bibliographystyle{alpha}
\bibliography{references}
\end{document}
